% Options for packages loaded elsewhere
\PassOptionsToPackage{unicode}{hyperref}
\PassOptionsToPackage{hyphens}{url}
%
\documentclass[
]{article}
\usepackage{amsmath,amssymb}
\usepackage{iftex}
\ifPDFTeX
  \usepackage[T1]{fontenc}
  \usepackage[utf8]{inputenc}
  \usepackage{textcomp} % provide euro and other symbols
\else % if luatex or xetex
  \usepackage{unicode-math} % this also loads fontspec
  \defaultfontfeatures{Scale=MatchLowercase}
  \defaultfontfeatures[\rmfamily]{Ligatures=TeX,Scale=1}
\fi
\usepackage{lmodern}
\ifPDFTeX\else
  % xetex/luatex font selection
\fi
% Use upquote if available, for straight quotes in verbatim environments
\IfFileExists{upquote.sty}{\usepackage{upquote}}{}
\IfFileExists{microtype.sty}{% use microtype if available
  \usepackage[]{microtype}
  \UseMicrotypeSet[protrusion]{basicmath} % disable protrusion for tt fonts
}{}
\makeatletter
\@ifundefined{KOMAClassName}{% if non-KOMA class
  \IfFileExists{parskip.sty}{%
    \usepackage{parskip}
  }{% else
    \setlength{\parindent}{0pt}
    \setlength{\parskip}{6pt plus 2pt minus 1pt}}
}{% if KOMA class
  \KOMAoptions{parskip=half}}
\makeatother
\usepackage{xcolor}
\usepackage[margin=1in]{geometry}
\usepackage{color}
\usepackage{fancyvrb}
\newcommand{\VerbBar}{|}
\newcommand{\VERB}{\Verb[commandchars=\\\{\}]}
\DefineVerbatimEnvironment{Highlighting}{Verbatim}{commandchars=\\\{\}}
% Add ',fontsize=\small' for more characters per line
\usepackage{framed}
\definecolor{shadecolor}{RGB}{248,248,248}
\newenvironment{Shaded}{\begin{snugshade}}{\end{snugshade}}
\newcommand{\AlertTok}[1]{\textcolor[rgb]{0.94,0.16,0.16}{#1}}
\newcommand{\AnnotationTok}[1]{\textcolor[rgb]{0.56,0.35,0.01}{\textbf{\textit{#1}}}}
\newcommand{\AttributeTok}[1]{\textcolor[rgb]{0.13,0.29,0.53}{#1}}
\newcommand{\BaseNTok}[1]{\textcolor[rgb]{0.00,0.00,0.81}{#1}}
\newcommand{\BuiltInTok}[1]{#1}
\newcommand{\CharTok}[1]{\textcolor[rgb]{0.31,0.60,0.02}{#1}}
\newcommand{\CommentTok}[1]{\textcolor[rgb]{0.56,0.35,0.01}{\textit{#1}}}
\newcommand{\CommentVarTok}[1]{\textcolor[rgb]{0.56,0.35,0.01}{\textbf{\textit{#1}}}}
\newcommand{\ConstantTok}[1]{\textcolor[rgb]{0.56,0.35,0.01}{#1}}
\newcommand{\ControlFlowTok}[1]{\textcolor[rgb]{0.13,0.29,0.53}{\textbf{#1}}}
\newcommand{\DataTypeTok}[1]{\textcolor[rgb]{0.13,0.29,0.53}{#1}}
\newcommand{\DecValTok}[1]{\textcolor[rgb]{0.00,0.00,0.81}{#1}}
\newcommand{\DocumentationTok}[1]{\textcolor[rgb]{0.56,0.35,0.01}{\textbf{\textit{#1}}}}
\newcommand{\ErrorTok}[1]{\textcolor[rgb]{0.64,0.00,0.00}{\textbf{#1}}}
\newcommand{\ExtensionTok}[1]{#1}
\newcommand{\FloatTok}[1]{\textcolor[rgb]{0.00,0.00,0.81}{#1}}
\newcommand{\FunctionTok}[1]{\textcolor[rgb]{0.13,0.29,0.53}{\textbf{#1}}}
\newcommand{\ImportTok}[1]{#1}
\newcommand{\InformationTok}[1]{\textcolor[rgb]{0.56,0.35,0.01}{\textbf{\textit{#1}}}}
\newcommand{\KeywordTok}[1]{\textcolor[rgb]{0.13,0.29,0.53}{\textbf{#1}}}
\newcommand{\NormalTok}[1]{#1}
\newcommand{\OperatorTok}[1]{\textcolor[rgb]{0.81,0.36,0.00}{\textbf{#1}}}
\newcommand{\OtherTok}[1]{\textcolor[rgb]{0.56,0.35,0.01}{#1}}
\newcommand{\PreprocessorTok}[1]{\textcolor[rgb]{0.56,0.35,0.01}{\textit{#1}}}
\newcommand{\RegionMarkerTok}[1]{#1}
\newcommand{\SpecialCharTok}[1]{\textcolor[rgb]{0.81,0.36,0.00}{\textbf{#1}}}
\newcommand{\SpecialStringTok}[1]{\textcolor[rgb]{0.31,0.60,0.02}{#1}}
\newcommand{\StringTok}[1]{\textcolor[rgb]{0.31,0.60,0.02}{#1}}
\newcommand{\VariableTok}[1]{\textcolor[rgb]{0.00,0.00,0.00}{#1}}
\newcommand{\VerbatimStringTok}[1]{\textcolor[rgb]{0.31,0.60,0.02}{#1}}
\newcommand{\WarningTok}[1]{\textcolor[rgb]{0.56,0.35,0.01}{\textbf{\textit{#1}}}}
\usepackage{graphicx}
\makeatletter
\newsavebox\pandoc@box
\newcommand*\pandocbounded[1]{% scales image to fit in text height/width
  \sbox\pandoc@box{#1}%
  \Gscale@div\@tempa{\textheight}{\dimexpr\ht\pandoc@box+\dp\pandoc@box\relax}%
  \Gscale@div\@tempb{\linewidth}{\wd\pandoc@box}%
  \ifdim\@tempb\p@<\@tempa\p@\let\@tempa\@tempb\fi% select the smaller of both
  \ifdim\@tempa\p@<\p@\scalebox{\@tempa}{\usebox\pandoc@box}%
  \else\usebox{\pandoc@box}%
  \fi%
}
% Set default figure placement to htbp
\def\fps@figure{htbp}
\makeatother
\setlength{\emergencystretch}{3em} % prevent overfull lines
\providecommand{\tightlist}{%
  \setlength{\itemsep}{0pt}\setlength{\parskip}{0pt}}
\setcounter{secnumdepth}{-\maxdimen} % remove section numbering
\usepackage{bookmark}
\IfFileExists{xurl.sty}{\usepackage{xurl}}{} % add URL line breaks if available
\urlstyle{same}
\hypersetup{
  pdftitle={Assignment 1},
  pdfauthor={Megan Hessel},
  hidelinks,
  pdfcreator={LaTeX via pandoc}}

\title{Assignment 1}
\usepackage{etoolbox}
\makeatletter
\providecommand{\subtitle}[1]{% add subtitle to \maketitle
  \apptocmd{\@title}{\par {\large #1 \par}}{}{}
}
\makeatother
\subtitle{California Spiny Lobster (\emph{Panulirus Interruptus}):
Assessing the Impact of Marine Protected Areas (MPAs) at 5 Reef Sites in
Santa Barbara County}
\author{Megan Hessel}
\date{1/8/26}

\begin{document}
\maketitle

\begin{center}\rule{0.5\linewidth}{0.5pt}\end{center}

\pandocbounded{\includegraphics[keepaspectratio]{figures/spiny2.jpg}}

\begin{center}\rule{0.5\linewidth}{0.5pt}\end{center}

\subsubsection{Assignment Instructions:}\label{assignment-instructions}

\begin{itemize}
\item
  Working with partners to troubleshoot code and concepts is encouraged!
  If you work with a partner, please list their name next to yours at
  the top of your assignment so Annie and I can easily see who
  collaborated.
\item
  All written responses must be written independently (\textbf{in your
  own words}).
\item
  Please follow the question prompts carefully and include only the
  information each question asks in your submitted responses.
\item
  Submit both your knitted document and the associated
  \texttt{RMarkdown} or \texttt{Quarto} file.
\item
  Your knitted presentation should meet the quality you'd submit to
  research colleagues or feel confident sharing publicly. Refer to the
  rubric for details about presentation standards.
\end{itemize}

\textbf{Assignment submission:} Megan Hessel

\begin{center}\rule{0.5\linewidth}{0.5pt}\end{center}

\begin{Shaded}
\begin{Highlighting}[]
\FunctionTok{library}\NormalTok{(tidyverse)}
\FunctionTok{library}\NormalTok{(here)}
\FunctionTok{library}\NormalTok{(janitor)}
\FunctionTok{library}\NormalTok{(estimatr)  }
\FunctionTok{library}\NormalTok{(performance)}
\FunctionTok{library}\NormalTok{(jtools)}
\FunctionTok{library}\NormalTok{(gt)}
\FunctionTok{library}\NormalTok{(gtsummary)}
\FunctionTok{library}\NormalTok{(interactions) }
\FunctionTok{library}\NormalTok{(dplyr)}
\FunctionTok{library}\NormalTok{(ggridges)}
\FunctionTok{library}\NormalTok{(paletteer)}
\FunctionTok{library}\NormalTok{(gtsummary)}
\FunctionTok{library}\NormalTok{(gt)}
\FunctionTok{library}\NormalTok{(MASS)}
\end{Highlighting}
\end{Shaded}

\begin{center}\rule{0.5\linewidth}{0.5pt}\end{center}

\paragraph{DATA SOURCE:}\label{data-source}

\begin{quote}
\href{https://doi.org/10.6073/pasta/a593a675d644fdefb736750b291579a0}{Reed
D. 2019. SBC LTER: Reef: Abundance, size and fishing effort for
California Spiny Lobster (Panulirus interruptus), ongoing since 2012.
Environmental Data Initiative.} Data accessed 11/17/2019.
\end{quote}

\begin{center}\rule{0.5\linewidth}{0.5pt}\end{center}

\subsubsection{\texorpdfstring{\textbf{Introduction}}{Introduction}}\label{introduction}

You're about to dive into some deep data collected from five reef sites
in Santa Barbara County, all about the abundance of California spiny
lobsters! 🦞 Data was gathered by divers annually from 2012 to 2018
across Naples, Mohawk, Isla Vista, Carpinteria, and Arroyo Quemado
reefs.

Why lobsters? Well, this sample provides an opportunity to evaluate the
impact of Marine Protected Areas (MPAs) established on January 1, 2012
(Reed, 2019). Of these five reefs, \emph{Naples, and Isla Vista are
MPAs,} while the other three are not protected (non-MPAs). Comparing
lobster health between these protected and non-protected areas gives us
the chance to study how commercial and recreational fishing might impact
these ecosystems.

We will consider the MPA sites the \texttt{treatment} group and use
regression methods to explore whether protecting these reefs really
makes a difference compared to non-MPA sites (our control group). In
this assignment, we'll think deeply about which causal inference
assumptions hold up under the research design and identify where they
fall short.

Let's break it down step by step and see what the data reveals! 📊

\pandocbounded{\includegraphics[keepaspectratio]{figures/map-5reefs.png}}

\begin{center}\rule{0.5\linewidth}{0.5pt}\end{center}

\paragraph{Step 1: Anticipating potential sources of selection
bias}\label{step-1-anticipating-potential-sources-of-selection-bias}

\textbf{a.} Do the control sites (Arroyo Quemado, Carpenteria, and
Mohawk) provide a strong counterfactual for our treatment sites (Naples,
Isla Vista)? Write a paragraph making a case for why this comparison is
ceteris paribus or whether selection bias is likely (be specific!).

There might be some selection bias due to the sites' differing in
current intensity and patterns, anthropomorphic activity, and
temperature. Point Conception (which is just West of this map) is where
the cold, subarctic Northern California current and warm Southern
California current meet, creating upwhelling. Therefore, Eastern sites
might be warmer and facing more current strength than the Western sites
due to the mixing and upwhelling West of Santa Barbara. Additionally,
Naples, IV, and Monhawk are closer to the city; therefore, those sites
might face more pollution and disturbance. However, while sites differ
on several confounding variables, these sites selections are still
useful and close to \emph{ceters paribus}. Control groups are both east
and west of the treatment group, minimizing the effect of the different
current and temperature pressures. All sites are very close to Santa
Barbara. Due to aquatic systems spatial autocorrelation, pollution in
one location will most likely be in other locations. Therefore, in this
example, there might be some selection bias, but it is pretty close to
\emph{celeris paribus} as locations face similar conditions except MPA
status.

\begin{center}\rule{0.5\linewidth}{0.5pt}\end{center}

\paragraph{Step 2: Read \& wrangle data}\label{step-2-read-wrangle-data}

\textbf{a.} Read in the raw data from the ``data'' folder named
\texttt{spiny\_abundance\_sb\_18.csv}. Name the data.frame
\texttt{rawdata}

\textbf{b.} Use the function \texttt{clean\_names()} from the
\texttt{janitor} package

\begin{Shaded}
\begin{Highlighting}[]
\CommentTok{\# HINT: check for coding of missing values (\textasciigrave{}na = "{-}99999"\textasciigrave{})}

\NormalTok{rawdata }\OtherTok{\textless{}{-}} \FunctionTok{read\_csv}\NormalTok{(here}\SpecialCharTok{::}\FunctionTok{here}\NormalTok{(}\StringTok{"data"}\NormalTok{, }\StringTok{"spiny\_abundance\_sb\_18.csv"}\NormalTok{), }
                    \AttributeTok{na =} \StringTok{\textquotesingle{}{-}99999\textquotesingle{}}\NormalTok{) }\SpecialCharTok{\%\textgreater{}\%} \CommentTok{\# converting \textquotesingle{}{-}9999\textquotesingle{} values to NA }
\NormalTok{    janitor}\SpecialCharTok{::}\FunctionTok{clean\_names}\NormalTok{()}
\end{Highlighting}
\end{Shaded}

\textbf{c.} Create a new \texttt{df} named \texttt{tidyata}. Using the
variable \texttt{site} (reef location) create a new variable
\texttt{reef} as a \texttt{factor} and add the following labels in the
order listed (i.e., re-order the \texttt{levels}):

\begin{verbatim}
"Arroyo Quemado", "Carpenteria", "Mohawk", "Isla Vista",  "Naples"
\end{verbatim}

\begin{Shaded}
\begin{Highlighting}[]
\CommentTok{\# Creating the ordered wanted }
\NormalTok{custom\_order }\OtherTok{\textless{}{-}} \FunctionTok{c}\NormalTok{(}\StringTok{"Arroyo Quemado"}\NormalTok{, }\StringTok{"Carpenteria"}\NormalTok{, }\StringTok{"Mohawk"}\NormalTok{, }\StringTok{"Isla Vista"}\NormalTok{, }\StringTok{"Naples"}\NormalTok{)}

\CommentTok{\# Making a \textasciigrave{}reef\textasciigrave{} a factor from \textasciigrave{}site\textasciigrave{} col}
\NormalTok{tidydata }\OtherTok{\textless{}{-}}\NormalTok{ rawdata }\SpecialCharTok{\%\textgreater{}\%} 
    \FunctionTok{mutate}\NormalTok{(}\AttributeTok{reef =} \FunctionTok{factor}\NormalTok{(rawdata}\SpecialCharTok{$}\NormalTok{site, }\AttributeTok{levels =} \FunctionTok{c}\NormalTok{(}\StringTok{"AQUE"}\NormalTok{, }\StringTok{"CARP"}\NormalTok{, }\StringTok{"MOHK"}\NormalTok{, }\StringTok{"IVEE"}\NormalTok{,}\StringTok{"NAPL"}\NormalTok{), }\AttributeTok{labels =}\NormalTok{ custom\_order))}


\CommentTok{\# Checking the level order }
\FunctionTok{levels}\NormalTok{(tidydata}\SpecialCharTok{$}\NormalTok{reef)}
\end{Highlighting}
\end{Shaded}

Create new \texttt{df} named \texttt{spiny\_counts}

\textbf{d.} Create a new variable \texttt{counts} to allow for an
analysis of lobster counts where the unit-level of observation is the
total number of observed lobsters per \texttt{site}, \texttt{year} and
\texttt{transect}.

\begin{itemize}
\tightlist
\item
  Create a variable \texttt{mean\_size} from the variable
  \texttt{size\_mm}
\item
  NOTE: The variable \texttt{counts} should have values which are
  integers (whole numbers).
\item
  Make sure to account for missing cases (\texttt{na})!
\end{itemize}

\begin{Shaded}
\begin{Highlighting}[]
\CommentTok{\#HINT(d): Use \textasciigrave{}group\_by()\textasciigrave{} \& \textasciigrave{}summarize()\textasciigrave{} to provide the total number of lobsters observed at each site{-}year{-}transect row{-}observation. }

\NormalTok{spiny\_counts }\OtherTok{\textless{}{-}}\NormalTok{ tidydata }\SpecialCharTok{\%\textgreater{}\%} 
    \FunctionTok{group\_by}\NormalTok{(site, year, transect) }\SpecialCharTok{\%\textgreater{}\%} 
    \FunctionTok{summarise}\NormalTok{(}\AttributeTok{obs\_lob =} \FunctionTok{sum}\NormalTok{(count, }\AttributeTok{na.rm =} \ConstantTok{TRUE}\NormalTok{), }\CommentTok{\# Number of lobsters observed }
              \AttributeTok{mean\_size =} \FunctionTok{mean}\NormalTok{(size\_mm, }\AttributeTok{na.rm =} \ConstantTok{TRUE}\NormalTok{)) }\SpecialCharTok{\%\textgreater{}\%} 
    \FunctionTok{ungroup}\NormalTok{()}

\CommentTok{\# \textasciigrave{}count\textasciigrave{} = Total number of lobsters of this size in the lobster transect   }
\CommentTok{\# \textasciigrave{}num\_ao\textasciigrave{} = Number of lobster that only its antenna was visible}
\end{Highlighting}
\end{Shaded}

\textbf{e.} Create a new variable \texttt{mpa} with levels \texttt{MPA}
and \texttt{non\_MPA}. For our regression analysis create a numerical
variable \texttt{treat} where MPA sites are coded \texttt{1} and
non\_MPA sites are coded \texttt{0}

\begin{Shaded}
\begin{Highlighting}[]
\CommentTok{\#HINT(e): Use \textasciigrave{}case\_when()\textasciigrave{} to create the 3 new variable columns}

\NormalTok{spiny\_counts }\OtherTok{\textless{}{-}}\NormalTok{ spiny\_counts }\SpecialCharTok{\%\textgreater{}\%} 
    \FunctionTok{mutate}\NormalTok{(}\AttributeTok{mpa =} \FunctionTok{case\_when}\NormalTok{(}
\NormalTok{        site }\SpecialCharTok{==} \StringTok{\textquotesingle{}AQUE\textquotesingle{}} \SpecialCharTok{|}\NormalTok{ site }\SpecialCharTok{==} \StringTok{\textquotesingle{}CARP\textquotesingle{}} \SpecialCharTok{|}\NormalTok{ site }\SpecialCharTok{==} \StringTok{\textquotesingle{}MOHK\textquotesingle{}}  \SpecialCharTok{\textasciitilde{}} \StringTok{\textquotesingle{}non\_MPA\textquotesingle{}}\NormalTok{, }\CommentTok{\# non MPA sites }
\NormalTok{        site }\SpecialCharTok{==} \StringTok{\textquotesingle{}IVEE\textquotesingle{}} \SpecialCharTok{|}\NormalTok{ site }\SpecialCharTok{==} \StringTok{\textquotesingle{}NAPL\textquotesingle{}} \SpecialCharTok{\textasciitilde{}} \StringTok{\textquotesingle{}MPA\textquotesingle{}} \CommentTok{\# MPA sites }
\NormalTok{))}

\CommentTok{\# Making \textasciigrave{}mpa\textasciigrave{} and \textasciigrave{}treat\textasciigrave{} a factor with 0 / non\_MPA and 1 / MPA}
\NormalTok{custom\_order }\OtherTok{\textless{}{-}} \FunctionTok{c}\NormalTok{(}\StringTok{"non\_MPA"}\NormalTok{, }\StringTok{"MPA"}\NormalTok{)}
\CommentTok{\#spiny\_counts$mpa \textless{}{-} factor(spiny\_counts$mpa, levels = custom\_order)}
\NormalTok{spiny\_counts}\SpecialCharTok{$}\NormalTok{treat }\OtherTok{\textless{}{-}} \FunctionTok{factor}\NormalTok{(spiny\_counts}\SpecialCharTok{$}\NormalTok{mpa, }\AttributeTok{levels =}\NormalTok{ custom\_order, }\AttributeTok{labels =} \FunctionTok{c}\NormalTok{(}\DecValTok{0}\NormalTok{, }\DecValTok{1}\NormalTok{))}

\CommentTok{\# Checking levels of factors }
\CommentTok{\#levels(spiny\_counts$mpa)}
\FunctionTok{levels}\NormalTok{(spiny\_counts}\SpecialCharTok{$}\NormalTok{treat)}
\end{Highlighting}
\end{Shaded}

\begin{quote}
NOTE: This step is crucial to the analysis. Check with a friend or come
to TA/instructor office hours to make sure the counts are coded
correctly!
\end{quote}

\begin{center}\rule{0.5\linewidth}{0.5pt}\end{center}

\paragraph{Step 3: Explore \& visualize
data}\label{step-3-explore-visualize-data}

\textbf{a.} Take a look at the data! Get familiar with the data in each
\texttt{df} format (\texttt{tidydata}, \texttt{spiny\_counts})

\begin{Shaded}
\begin{Highlighting}[]
\FunctionTok{str}\NormalTok{(tidydata)}
\FunctionTok{str}\NormalTok{(spiny\_counts)}
\end{Highlighting}
\end{Shaded}

\textbf{b.} We will focus on the variables \texttt{count},
\texttt{year}, \texttt{site}, and \texttt{treat}(\texttt{mpa}) to model
lobster abundance. Create the following 4 plots using a different method
each time from the 6 options provided. Add a layer (\texttt{geom}) to
each of the plots including informative descriptive statistics (you
choose; e.g., mean, median, SD, quartiles, range). Make sure each plot
dimension is clearly labeled (e.g., axes, groups).

\begin{itemize}
\tightlist
\item
  \href{https://r-charts.com/distribution/density-plot-group-ggplot2}{Density
  plot}
\item
  \href{https://r-charts.com/distribution/ggridges/}{Ridge plot}
\item
  \href{https://ggplot2.tidyverse.org/reference/geom_jitter.html}{Jitter
  plot}
\item
  \href{https://r-charts.com/distribution/violin-plot-group-ggplot2}{Violin
  plot}
\item
  \href{https://r-charts.com/distribution/histogram-density-ggplot2/}{Histogram}
\item
  \href{https://r-charts.com/distribution/beeswarm/}{Beeswarm}
\end{itemize}

Create plots displaying the distribution of lobster \textbf{counts}:

\begin{enumerate}
\def\labelenumi{\arabic{enumi})}
\tightlist
\item
  grouped by reef site\\
\end{enumerate}

\begin{Shaded}
\begin{Highlighting}[]
\NormalTok{color\_palette }\OtherTok{\textless{}{-}} \FunctionTok{c}\NormalTok{(}\StringTok{\textquotesingle{}\#0a6664\textquotesingle{}}\NormalTok{, }\StringTok{\textquotesingle{}\#17BEBBFF\textquotesingle{}}\NormalTok{, }\StringTok{\textquotesingle{}\#D4F4DDFF\textquotesingle{}}\NormalTok{, }\StringTok{\textquotesingle{}\#D62246FF\textquotesingle{}}\NormalTok{, }\StringTok{\textquotesingle{}\#4B1D3FFF\textquotesingle{}}\NormalTok{)}

\CommentTok{\# Density graph by reef site }
\NormalTok{spiny\_counts }\SpecialCharTok{\%\textgreater{}\%} 
    \FunctionTok{mutate}\NormalTok{(}\AttributeTok{site =} \FunctionTok{factor}\NormalTok{(site, }
                         \AttributeTok{levels =} \FunctionTok{c}\NormalTok{(}\StringTok{"AQUE"}\NormalTok{, }\StringTok{"CARP"}\NormalTok{, }\StringTok{"MOHK"}\NormalTok{, }\StringTok{"IVEE"}\NormalTok{,}\StringTok{"NAPL"}\NormalTok{) }
\NormalTok{                         )) }\SpecialCharTok{\%\textgreater{}\%} 
    \FunctionTok{ggplot}\NormalTok{(}\FunctionTok{aes}\NormalTok{(}\AttributeTok{x =}\NormalTok{ obs\_lob, }\AttributeTok{y =}\NormalTok{ site, }\AttributeTok{fill =}\NormalTok{ site)) }\SpecialCharTok{+} 
    \FunctionTok{geom\_density\_ridges}\NormalTok{( }\AttributeTok{alpha =} \FloatTok{0.5}\NormalTok{, }
                         \AttributeTok{quantile\_lines =} \ConstantTok{TRUE}\NormalTok{) }\SpecialCharTok{+} 
    \FunctionTok{theme\_gray}\NormalTok{() }\SpecialCharTok{+} 
    \FunctionTok{labs}\NormalTok{(}\AttributeTok{x =} \StringTok{"Lobster Counts"}\NormalTok{, }
         \AttributeTok{y =} \StringTok{"Sites"}\NormalTok{,}
         \AttributeTok{title =} \StringTok{"California Spiny Lobster Abundance at Sample Sites"}\NormalTok{, }
         \AttributeTok{caption =} \StringTok{"Source: Reed, 2019"}\NormalTok{) }\SpecialCharTok{+} 
    \CommentTok{\#scale\_fill\_manual(values = color\_palette) + }
    \FunctionTok{theme}\NormalTok{(}\AttributeTok{legend.position=}\StringTok{"none"}\NormalTok{)}
\end{Highlighting}
\end{Shaded}

\begin{enumerate}
\def\labelenumi{\arabic{enumi})}
\setcounter{enumi}{1}
\tightlist
\item
  grouped by MPA status
\end{enumerate}

\begin{Shaded}
\begin{Highlighting}[]
\FunctionTok{ggplot}\NormalTok{(spiny\_counts, }\FunctionTok{aes}\NormalTok{(}\AttributeTok{x =}\NormalTok{ mpa, }\AttributeTok{y =}\NormalTok{ obs\_lob, }\AttributeTok{fill =}\NormalTok{ mpa)) }\SpecialCharTok{+} 
    \FunctionTok{geom\_violin}\NormalTok{(}\AttributeTok{alpha =} \FloatTok{0.5}\NormalTok{) }\SpecialCharTok{+} 
    \FunctionTok{scale\_fill\_manual}\NormalTok{(}\AttributeTok{values =} \FunctionTok{c}\NormalTok{(}\StringTok{\textquotesingle{}\#488f30\textquotesingle{}}\NormalTok{, }\StringTok{\textquotesingle{}\#ffad76\textquotesingle{}}\NormalTok{)) }\SpecialCharTok{+} 
    \FunctionTok{theme\_light}\NormalTok{() }\SpecialCharTok{+} 
    \FunctionTok{theme}\NormalTok{(}\AttributeTok{legend.position=}\StringTok{"none"}\NormalTok{) }\SpecialCharTok{+} 
    \FunctionTok{labs}\NormalTok{(}\AttributeTok{x =} \StringTok{\textquotesingle{}MPA Status\textquotesingle{}}\NormalTok{, }
         \AttributeTok{y =} \StringTok{\textquotesingle{}Lobster Count\textquotesingle{}}\NormalTok{, }
         \AttributeTok{title =} \StringTok{\textquotesingle{}California Spiny Lobster Abundance at MPA vs non{-}MPA Sites\textquotesingle{}}\NormalTok{, }
         \AttributeTok{caption =} \StringTok{"Source: Reed, 2019"}\NormalTok{) }\SpecialCharTok{+} 
    \FunctionTok{geom\_boxplot}\NormalTok{(}\AttributeTok{width =} \FloatTok{0.06}\NormalTok{, }
                 \AttributeTok{outlier.size =} \FloatTok{0.55}\NormalTok{, }
                 \AttributeTok{fill =} \StringTok{"grey90"}\NormalTok{)}
\end{Highlighting}
\end{Shaded}

\begin{enumerate}
\def\labelenumi{\arabic{enumi})}
\setcounter{enumi}{2}
\tightlist
\item
  grouped by year
\end{enumerate}

\begin{Shaded}
\begin{Highlighting}[]
\NormalTok{spiny\_summary }\OtherTok{\textless{}{-}}\NormalTok{ spiny\_counts }\SpecialCharTok{\%\textgreater{}\%} 
    \FunctionTok{group\_by}\NormalTok{(year) }\SpecialCharTok{\%\textgreater{}\%} 
    \FunctionTok{summarise}\NormalTok{(}\AttributeTok{mean\_count =} \FunctionTok{mean}\NormalTok{(obs\_lob))}


\FunctionTok{ggplot}\NormalTok{() }\SpecialCharTok{+} 
  \FunctionTok{geom\_jitter}\NormalTok{(}\AttributeTok{data =}\NormalTok{ spiny\_counts,}
    \AttributeTok{mapping =} \FunctionTok{aes}\NormalTok{(}\AttributeTok{x =} \FunctionTok{factor}\NormalTok{(year), }\AttributeTok{y =}\NormalTok{ obs\_lob,}
        \AttributeTok{color =} \StringTok{"Observations"}\NormalTok{),}
    \AttributeTok{width =} \FloatTok{0.25}\NormalTok{) }\SpecialCharTok{+} 
  \FunctionTok{geom\_jitter}\NormalTok{(}\AttributeTok{data =}\NormalTok{ spiny\_summary,}
    \AttributeTok{mapping =} \FunctionTok{aes}\NormalTok{(}\AttributeTok{x =} \FunctionTok{factor}\NormalTok{(year), }\AttributeTok{y =}\NormalTok{ mean\_count,}
        \AttributeTok{color =} \StringTok{"Mean abundance"}\NormalTok{),}
    \AttributeTok{width =} \FloatTok{0.25}\NormalTok{,}
    \AttributeTok{size =} \FloatTok{1.7}\NormalTok{) }\SpecialCharTok{+} 
  \FunctionTok{scale\_color\_manual}\NormalTok{(}
    \AttributeTok{name =} \StringTok{""}\NormalTok{,}
    \AttributeTok{values =} \FunctionTok{c}\NormalTok{(}
      \StringTok{"Observations"} \OtherTok{=} \StringTok{"grey15"}\NormalTok{,}
      \StringTok{"Mean abundance"} \OtherTok{=} \StringTok{"red2"}\NormalTok{)) }\SpecialCharTok{+}
  \FunctionTok{labs}\NormalTok{(}
    \AttributeTok{x =} \StringTok{"Year"}\NormalTok{,}
    \AttributeTok{y =} \StringTok{"Lobster Abundance"}\NormalTok{,}
    \AttributeTok{title =} \StringTok{"California Spiny Lobster Abundance Over Time"}\NormalTok{,}
    \AttributeTok{caption =} \StringTok{"Source: Reed, 2019"}\NormalTok{) }\SpecialCharTok{+}
  \FunctionTok{theme\_grey}\NormalTok{()}
\end{Highlighting}
\end{Shaded}

Create a plot of lobster \textbf{size} :

\begin{enumerate}
\def\labelenumi{\arabic{enumi})}
\setcounter{enumi}{3}
\tightlist
\item
  You choose the grouping variable(s)!
\end{enumerate}

\begin{Shaded}
\begin{Highlighting}[]
\CommentTok{\# Calculate median }
\NormalTok{median\_group }\OtherTok{\textless{}{-}}\NormalTok{ spiny\_counts }\SpecialCharTok{\%\textgreater{}\%} 
    \FunctionTok{group\_by}\NormalTok{(mpa) }\SpecialCharTok{\%\textgreater{}\%} 
    \FunctionTok{summarise}\NormalTok{(}\AttributeTok{median =} \FunctionTok{median}\NormalTok{(mean\_size, }\AttributeTok{na.rm =} \ConstantTok{TRUE}\NormalTok{, }\AttributeTok{.group =}\NormalTok{ drop))}

\CommentTok{\# Plot with medians }
\FunctionTok{ggplot}\NormalTok{() }\SpecialCharTok{+} 
    \FunctionTok{geom\_density}\NormalTok{(}
        \AttributeTok{data =}\NormalTok{ spiny\_counts,}
        \AttributeTok{mapping =} \FunctionTok{aes}\NormalTok{(}\AttributeTok{x =}\NormalTok{ mean\_size, }\AttributeTok{fill =}\NormalTok{ mpa),}
        \AttributeTok{alpha =} \FloatTok{0.4}\NormalTok{) }\SpecialCharTok{+} 
    \FunctionTok{theme\_minimal}\NormalTok{() }\SpecialCharTok{+}
    \FunctionTok{geom\_vline}\NormalTok{(}\AttributeTok{data =}\NormalTok{ median\_group,}
        \AttributeTok{mapping =}\NormalTok{ (}\FunctionTok{aes}\NormalTok{(}\AttributeTok{xintercept =}\NormalTok{ median, }
                       \AttributeTok{color =}\NormalTok{ mpa)), }
        \AttributeTok{linetype =} \StringTok{"dashed"}\NormalTok{) }\SpecialCharTok{+} 
    \FunctionTok{scale\_color\_manual}\NormalTok{(}\AttributeTok{values =} \FunctionTok{c}\NormalTok{(}\StringTok{"darkgreen"}\NormalTok{, }\StringTok{"darkorange3"}\NormalTok{))}\SpecialCharTok{+}
    \FunctionTok{scale\_fill\_manual}\NormalTok{(}\AttributeTok{values =} \FunctionTok{c}\NormalTok{(}\StringTok{"\#488f30"}\NormalTok{, }\StringTok{"\#ffad76"}\NormalTok{)) }\SpecialCharTok{+} 
    \FunctionTok{labs}\NormalTok{(}\AttributeTok{x =} \StringTok{"Lobster Mean Size (mm)"}\NormalTok{, }
         \AttributeTok{y =} \StringTok{"Density"}\NormalTok{, }
         \AttributeTok{title =} \StringTok{"California Spiny Lobster Size in MPA and non{-}MPA Sites"}\NormalTok{, }
         \AttributeTok{fill =} \StringTok{"MPA Status"}\NormalTok{, }
         \AttributeTok{color =} \StringTok{"Size Median Values"}\NormalTok{) }
\end{Highlighting}
\end{Shaded}

\textbf{c.} Compare means of the outcome by treatment group. Using the
\texttt{tbl\_summary()} function from the package
\href{https://www.danieldsjoberg.com/gtsummary/articles/tbl_summary.html}{\texttt{gt\_summary}}

\begin{Shaded}
\begin{Highlighting}[]
\NormalTok{spiny\_counts }\SpecialCharTok{\%\textgreater{}\%} 
    \FunctionTok{rename}\NormalTok{(}\AttributeTok{lobster\_count =}\NormalTok{ obs\_lob) }\SpecialCharTok{\%\textgreater{}\%} 
\NormalTok{    gtsummary}\SpecialCharTok{::}\FunctionTok{tbl\_summary}\NormalTok{(}
        \AttributeTok{by =}\NormalTok{ treat, }
        \AttributeTok{include =} \FunctionTok{c}\NormalTok{(lobster\_count, mean\_size, site), }
        \AttributeTok{statistic =} \FunctionTok{list}\NormalTok{(}\FunctionTok{all\_continuous}\NormalTok{() }\SpecialCharTok{\textasciitilde{}} \StringTok{\textquotesingle{}\{mean\} (\{sd\})\textquotesingle{}}\NormalTok{)}
\NormalTok{    ) }\SpecialCharTok{\%\textgreater{}\%} 
    \FunctionTok{add\_p}\NormalTok{() }\SpecialCharTok{\%\textgreater{}\%} 
    \CommentTok{\# Add header labels }
    \FunctionTok{modify\_header}\NormalTok{(label }\SpecialCharTok{\textasciitilde{}} \StringTok{"**Variable**"}\NormalTok{) }\SpecialCharTok{\%\textgreater{}\%}
    \FunctionTok{modify\_spanning\_header}\NormalTok{(}\FunctionTok{c}\NormalTok{(}\StringTok{"stat\_1"}\NormalTok{, }\StringTok{"stat\_2"}\NormalTok{) }\SpecialCharTok{\textasciitilde{}} \StringTok{"**Treatment**"}\NormalTok{) }
\end{Highlighting}
\end{Shaded}

\begin{center}\rule{0.5\linewidth}{0.5pt}\end{center}

\paragraph{Step 4: OLS regression- building
intuition}\label{step-4-ols-regression--building-intuition}

\textbf{a.} Start with a simple OLS estimator of lobster counts
regressed on treatment. Use the function \texttt{summ()} from the
\href{https://jtools.jacob-long.com/}{\texttt{jtools}} package to print
the OLS output

\textbf{b.} Interpret the intercept \& predictor coefficients \emph{in
your own words}. Use full sentences and write your interpretation of the
regression results to be as clear as possible to a non-academic
audience.

\begin{Shaded}
\begin{Highlighting}[]
\CommentTok{\# }\AlertTok{NOTE}\CommentTok{: We will not evaluate/interpret model fit in this assignment (e.g., R{-}square)}

\NormalTok{m1\_ols }\OtherTok{\textless{}{-}} \FunctionTok{lm}\NormalTok{(obs\_lob }\SpecialCharTok{\textasciitilde{}}\NormalTok{ treat, }\AttributeTok{data =}\NormalTok{ spiny\_counts)}

\FunctionTok{summ}\NormalTok{(m1\_ols, }\AttributeTok{model.fit =} \ConstantTok{FALSE}\NormalTok{) }
\end{Highlighting}
\end{Shaded}

\emph{Intercept:} With no treatment (in a non-MPA site), there is
expected to be an average of \textasciitilde22.75 lobsters.

\emph{Coefficient Estimation:} With treatment (MPA protection), lobsters
count is expected to increase by approximately 5.36 lobsters.

\textbf{c.} Check the model assumptions using the \texttt{check\_model}
function from the \texttt{performance} package

\textbf{d.} Explain the results of the 4 diagnostic plots. Why are we
getting this result?

\begin{Shaded}
\begin{Highlighting}[]
\FunctionTok{check\_model}\NormalTok{(m1\_ols,  }\AttributeTok{check =} \StringTok{"qq"}\NormalTok{ )}
\end{Highlighting}
\end{Shaded}

\begin{Shaded}
\begin{Highlighting}[]
\FunctionTok{check\_model}\NormalTok{(m1\_ols, }\AttributeTok{check =} \StringTok{"normality"}\NormalTok{)}
\end{Highlighting}
\end{Shaded}

\begin{Shaded}
\begin{Highlighting}[]
\FunctionTok{check\_model}\NormalTok{(m1\_ols, }\AttributeTok{check =} \StringTok{"homogeneity"}\NormalTok{)}
\end{Highlighting}
\end{Shaded}

\begin{Shaded}
\begin{Highlighting}[]
\FunctionTok{check\_model}\NormalTok{(m1\_ols, }\AttributeTok{check =} \StringTok{"pp\_check"}\NormalTok{)}
\end{Highlighting}
\end{Shaded}

By the \texttt{check\_model} graphs, the data is not normality
distributed (with a leptokurtic skewed distribution), the residuals are
not normality distribution, and there is no equal variance
(heteroscedasticity). The \texttt{check\_model} results are showing the
data's inherently skewed characteristic. These violate OSL model's
assumptions.

\begin{center}\rule{0.5\linewidth}{0.5pt}\end{center}

\paragraph{Step 5: Fitting GLMs}\label{step-5-fitting-glms}

\textbf{a.} Estimate a \emph{Poisson regression model} using the
\texttt{glm()} function

\textbf{b.} Interpret the predictor coefficient in your own words. Use
full sentences and write your interpretation of the results to be as
clear as possible to a non-academic audience.

\emph{Intercept:} exp(3.12) = 22.64638 \emph{Est coefficient:} exp(0.21)
= 1.233678 With no treatment (non-MPA locations), the lobster count is
expected to be \textasciitilde23 on average. With treatment (MPA
protection), lobster count is expected to increase by 23\% which is
statistically significant.

\[ 
# Yield Percent Change
(exp(estimation) - 1) * 100
\]

\textbf{c.} Explain the statistical concept of dispersion and
overdispersion in the context of this model.

Dispersion is when the mean and variance are similar in portion to each
other. Overdispersion is when there is greater variance than expected
and thus not proportional to the mean. Poisson is a statistical test
where mean and variance must be proportional; overdispersion in Poisson
can cause misleading test statistics and standard errors. A dispersion
ratio of 67.033 indicates overdispersion which violates the Poisson
assumption.

\textbf{d.} Compare results with previous model, explain change in the
significance of the treatment effect

Both models have similar predictions. However, in the OSL model, there
is no statistical significance in lobster counts between MPA and non-MPA
sites. In the Poisson model, there is a statistical significance in
lobster abundance between treatments. The Poisson model is better fitted
for count data, and therefore give more accurate results.

\textbf{e.} Check the model assumptions. Explain results.

This data have overdispersion (dispersion ration of 67.033) and
zero-inflation (27 observed zeros) which breaks Poisson model
assumptions. Therefore, the Poisson model is inappropriate for this data
set.

\textbf{f.} Conduct tests for over-dispersion \& zero-inflation. Explain
results.

\begin{Shaded}
\begin{Highlighting}[]
\CommentTok{\# Poisson model }

\NormalTok{m2\_pois }\OtherTok{\textless{}{-}} \FunctionTok{glm}\NormalTok{(obs\_lob }\SpecialCharTok{\textasciitilde{}}\NormalTok{ treat, }\AttributeTok{data =}\NormalTok{ spiny\_counts, }\AttributeTok{family =} \FunctionTok{poisson}\NormalTok{(}\AttributeTok{link =} \StringTok{"log"}\NormalTok{))}

\FunctionTok{summ}\NormalTok{(m2\_pois, }\AttributeTok{model.fit =} \ConstantTok{FALSE}\NormalTok{)}
\end{Highlighting}
\end{Shaded}

\begin{Shaded}
\begin{Highlighting}[]
\CommentTok{\# Checking assumptions}
\FunctionTok{check\_model}\NormalTok{(m2\_pois)}

\FunctionTok{check\_overdispersion}\NormalTok{(m2\_pois)}

\FunctionTok{check\_zeroinflation}\NormalTok{(m2\_pois)}
\end{Highlighting}
\end{Shaded}

\textbf{g.} Fit a \emph{negative binomial model} using the function
glm.nb() from the package \texttt{MASS} and check model diagnostics

\textbf{h.} In 1-2 sentences explain rationale for fitting this GLM
model.

Negative binomial model is more flexible than the Poisson model because
it accounts for overdispersion by adding a dispersion parameter. Because
this data has overdispersion, it is more appropriately fitted for the
negative binomial model.

\textbf{i.} Interpret the treatment estimate result in your own words.
Compare with results from the previous model.

exp(3.12) = 22.64638 exp(0.21) = 1.233678

In not protected areas, the expected lobster count is \textasciitilde23
on average. In MPA sites, there is a expected increase of 23\% lobsters,
but this relationships is not significant.

All models suggest lobster abundance is higher in Marine Protected Areas
(MPAs). Specifically, the estimations from the poisson and negative
binomial models are very similar. The P-values are the main difference
between the models. Poisson shows statistical significance while the
negative binomal model and linear regression does not. The Poission
P-value result should be treated with caution because the assumptions
for Poisson are violated which can lead to misleading P-values.
Similarly, the linear regession assumptions are violated which can cause
unreliable results.

\begin{Shaded}
\begin{Highlighting}[]
\CommentTok{\# Negative Binomal }
\NormalTok{m3\_nb }\OtherTok{\textless{}{-}} \FunctionTok{glm.nb}\NormalTok{(obs\_lob }\SpecialCharTok{\textasciitilde{}}\NormalTok{ treat, }\AttributeTok{data =}\NormalTok{ spiny\_counts) }\CommentTok{\# \textasciigrave{}glm.nb()\textasciigrave{}  does not require a \textasciigrave{}family\textasciigrave{} argument}

\FunctionTok{summ}\NormalTok{(m3\_nb, }\AttributeTok{model.fit =} \ConstantTok{FALSE}\NormalTok{)}
\end{Highlighting}
\end{Shaded}

\begin{Shaded}
\begin{Highlighting}[]
\FunctionTok{check\_overdispersion}\NormalTok{(m3\_nb)}

\FunctionTok{check\_zeroinflation}\NormalTok{(m3\_nb)}
\end{Highlighting}
\end{Shaded}

\begin{Shaded}
\begin{Highlighting}[]
\FunctionTok{check\_predictions}\NormalTok{(m3\_nb)}

\FunctionTok{check\_model}\NormalTok{(m3\_nb)}
\end{Highlighting}
\end{Shaded}

\begin{center}\rule{0.5\linewidth}{0.5pt}\end{center}

\paragraph{Step 6: Compare models}\label{step-6-compare-models}

\textbf{a.} Use the \texttt{export\_summ()} function from the
\texttt{jtools} package to look at the three regression models you fit
side-by-side.

\textbf{c.} Write a short paragraph comparing the results. Is the
treatment effect \texttt{robust} or stable across the model
specifications.

In all three models, the treatment coefficient is positive, meaning
there is an increases of lobster abundance in the treatment sites (MPAs)
compared to the untreated (unprotected) areas. Additionally, the
estimates themselves were all very similar between models. The primary
distinctions between the models are the P-values and standard
deviations. The OLS model shows high uncertainty in its positive
prediction and no statistical significance. The OLS is a poor model for
the count dataset, seen by the high uncertainty. The Poisson model's
P-value shows statistical significance. However, the over-dispersion and
zero-inflated nature of the count data makes Poisson an inadequate
model. Lastly, the negative binomial accounts for over-dispersion and
zero values, making the model a more appropriate fit. The negative
binomial model estimates a positive treatment effect, but no statistical
significance. In conclusion, the estimates are all consistent. However,
statistical significance is not robust across all the models
specifications.

\begin{Shaded}
\begin{Highlighting}[]
\FunctionTok{export\_summs}\NormalTok{(m1\_ols, m2\_pois, m3\_nb,}
             \AttributeTok{model.names =} \FunctionTok{c}\NormalTok{(}\StringTok{"OLS"}\NormalTok{,}\StringTok{"Poisson"}\NormalTok{, }\StringTok{"NB"}\NormalTok{),}
             \AttributeTok{statistics =} \StringTok{"none"}\NormalTok{)}
\end{Highlighting}
\end{Shaded}

\begin{center}\rule{0.5\linewidth}{0.5pt}\end{center}

\paragraph{Step 7: Building intuition - fixed
effects}\label{step-7-building-intuition---fixed-effects}

\textbf{a.} Create new \texttt{df} with the \texttt{year} variable
converted to a factor

\textbf{b.} Run the following negative binomial model using
\texttt{glm.nb()}

\begin{itemize}
\tightlist
\item
  Add fixed effects for \texttt{year} (i.e., dummy coefficients)
\item
  Include an interaction term between variables \texttt{treat} \&
  \texttt{year} (\texttt{treat*year})
\end{itemize}

\textbf{c.} Take a look at the regression output. Each coefficient
provides a comparison or the difference in means for a specific
sub-group in the data. Informally, describe the what the model has
estimated at a conceptual level (NOTE: you do not have to interpret
coefficients individually)

By adding \texttt{year} as an additive and multiplicative variable, we
are now examining the change of lobster counts in MPA vs non-MPA sites
\emph{overtime}. The reference year is 2012. The year coefficients are
the expected average change of lobster abundance from 2012 in
unprotected sites. The \texttt{treat} coefficient explains the expected
change of lobster abundance between treatment (MPAs) and non-treatment
(non-MPA) sites - the effect of the treatment. The interaction between
the predictors examines if lobster counts in MPAs are changing relative
to non-MPA sites in each year.

\textbf{d.} Explain why the main effect for treatment is negative? *Does
this result make sense?

The treatment coefficient examines the change of lobster counts when the
year is fixed (fixed to the reference year - 2012). The results show, in
2012, the MPA sites had on average 82\% less lobsters than the non-MPA
sites. MPA lobster abundance increases in the later years which is seen
in the interaction terms, not the treatment coefficient. These results
make sense as MPA are designed to increase biodiversity overtime.

\begin{Shaded}
\begin{Highlighting}[]
\NormalTok{ff\_counts }\OtherTok{\textless{}{-}}\NormalTok{ spiny\_counts }\SpecialCharTok{\%\textgreater{}\%} 
    \FunctionTok{mutate}\NormalTok{(}\AttributeTok{year=}\FunctionTok{as\_factor}\NormalTok{(year))}
    
\NormalTok{m5\_fixedeffs }\OtherTok{\textless{}{-}} \FunctionTok{glm.nb}\NormalTok{(}
\NormalTok{    obs\_lob }\SpecialCharTok{\textasciitilde{}} 
\NormalTok{        treat }\SpecialCharTok{+}
\NormalTok{        year }\SpecialCharTok{+}
\NormalTok{        treat}\SpecialCharTok{*}\NormalTok{year,}
    \AttributeTok{data =}\NormalTok{ ff\_counts)}

\FunctionTok{summ}\NormalTok{(m5\_fixedeffs, }\AttributeTok{model.fit =} \ConstantTok{FALSE}\NormalTok{)}
\end{Highlighting}
\end{Shaded}

\textbf{e.} Look at the model predictions: Use the
\texttt{interact\_plot()} function from package \texttt{interactions} to
plot mean predictions by year and treatment status.

\textbf{f.} Re-evaluate your responses (c) and (b) above.

My previous explanation of the terms is exemplified in the
\texttt{interact\_plot()}. In 2012 (the reference year), MPA sites had
considerably less lobster abundance. As time continued, lobster counts
increased overall, especially in treatment/protected areas (with a drop
in the population in 2016 at MPA sites).

\begin{Shaded}
\begin{Highlighting}[]
\CommentTok{\# What teacher did }
\FunctionTok{interact\_plot}\NormalTok{(}\AttributeTok{model =}\NormalTok{ m5\_fixedeffs, }
              \AttributeTok{pred =}\NormalTok{ year, }\CommentTok{\# Predictor variable}
              \AttributeTok{modx =}\NormalTok{ treat, }\CommentTok{\# Moderator variable}
              \AttributeTok{outcome.scale =} \StringTok{"link"}\NormalTok{) }\CommentTok{\# }\AlertTok{NOTE}\CommentTok{: y{-}axis on log{-}scale}

\CommentTok{\# HINT: Change \textasciigrave{}outcome.scale\textasciigrave{} to "response" to convert y{-}axis scale to counts}
\end{Highlighting}
\end{Shaded}

\textbf{g.} Using \texttt{ggplot()} create a plot in same style as the
previous \texttt{interaction\ plot}, but displaying the original scale
of the outcome variable (lobster counts). This type of plot is commonly
used to show how the treatment effect changes across discrete time
points (i.e., panel data).

\begin{Shaded}
\begin{Highlighting}[]
\CommentTok{\# Making the dataframe }
\NormalTok{plot\_counts }\OtherTok{\textless{}{-}}\NormalTok{ ff\_counts }\SpecialCharTok{\%\textgreater{}\%} 
    \FunctionTok{group\_by}\NormalTok{(year, treat) }\SpecialCharTok{\%\textgreater{}\%} 
    \FunctionTok{summarise}\NormalTok{(}\AttributeTok{mean\_count =} \FunctionTok{mean}\NormalTok{(obs\_lob)) }\SpecialCharTok{\%\textgreater{}\%} 
    \FunctionTok{ungroup}\NormalTok{()}

\CommentTok{\# Plotting }
\NormalTok{plot\_counts }\SpecialCharTok{\%\textgreater{}\%}
    \FunctionTok{mutate}\NormalTok{(}\AttributeTok{mpa =} \FunctionTok{recode}\NormalTok{(treat, }\CommentTok{\# Rename mpa characters }
                        \StringTok{\textasciigrave{}}\AttributeTok{0}\StringTok{\textasciigrave{}} \OtherTok{=} \StringTok{"non{-}MPA"}\NormalTok{, }
                        \StringTok{\textasciigrave{}}\AttributeTok{1}\StringTok{\textasciigrave{}} \OtherTok{=} \StringTok{"MPA"}\NormalTok{)) }\SpecialCharTok{\%\textgreater{}\%} 
    \FunctionTok{ggplot}\NormalTok{(}\FunctionTok{aes}\NormalTok{(}\AttributeTok{x =}\NormalTok{ year, }\AttributeTok{y =}\NormalTok{ mean\_count, }\AttributeTok{color =}\NormalTok{ mpa, }\AttributeTok{group =}\NormalTok{ mpa)) }\SpecialCharTok{+} \CommentTok{\# plot }
    \FunctionTok{geom\_line}\NormalTok{(}\AttributeTok{linewidth =} \DecValTok{1}\NormalTok{)}\SpecialCharTok{+} 
    \FunctionTok{geom\_point}\NormalTok{(}\AttributeTok{size =} \DecValTok{3}\NormalTok{) }\SpecialCharTok{+} 
    \FunctionTok{scale\_color\_manual}\NormalTok{(}\AttributeTok{values =} \FunctionTok{c}\NormalTok{(}\StringTok{"\#ffad76"}\NormalTok{, }\StringTok{"\#488f30"}\NormalTok{)) }\SpecialCharTok{+} 
    \FunctionTok{theme\_minimal}\NormalTok{() }\SpecialCharTok{+} 
    \FunctionTok{labs}\NormalTok{(}
        \AttributeTok{title =} \StringTok{"Lobster Abundance Overtime"}\NormalTok{, }
        \AttributeTok{x =} \StringTok{"Year"}\NormalTok{, }
        \AttributeTok{y =} \StringTok{"Mean Lobster Count"}\NormalTok{, }
        \AttributeTok{color =} \StringTok{"MPA Status"}
\NormalTok{    )}
\end{Highlighting}
\end{Shaded}

\begin{center}\rule{0.5\linewidth}{0.5pt}\end{center}

\paragraph{Step 8: Reconsider causal identification
assumptions}\label{step-8-reconsider-causal-identification-assumptions}

\begin{enumerate}
\def\labelenumi{\alph{enumi}.}
\tightlist
\item
  Discuss whether you think \texttt{spillover\ effects} are likely in
  this research context (see Glossary of terms;
  \url{https://docs.google.com/document/d/1RIudsVcYhWGpqC-Uftk9UTz3PIq6stVyEpT44EPNgpE/edit?usp=sharing})
\end{enumerate}

Spillover is when the treatment and control outcomes can impact each
other - outcomes are not independent of each other. In marine MPA
situations, spillover is a factor. Marine life does not know MPA extends
and therefore travel inside and outside the area, impacting lobster
count in treated and non-treated areas. Spillover effects are strongest
in the area surrounding the MPA. In this research, the non-MPA area are
far way from MPAs which makes the spillover effect less intense.

\begin{enumerate}
\def\labelenumi{\alph{enumi}.}
\setcounter{enumi}{1}
\tightlist
\item
  Explain why spillover is an issue for the identification of causal
  effects.
\end{enumerate}

Spillover effects make it hard to discover the causal impacts of the
treatment when the outcomes depend on each other. The intermixing of
outcomes makes it difficult to disguising the specific cause and effect.

\begin{enumerate}
\def\labelenumi{\alph{enumi}.}
\setcounter{enumi}{2}
\item
  Discuss the following causal inference assumptions in the context of
  the MPA treatment effect estimator. Evaluate if each of the assumption
  are reasonable:

  \begin{enumerate}
  \def\labelenumii{\arabic{enumii})}
  \tightlist
  \item
    SUTVA: Stable Unit Treatment Value assumption
  \item
    `Exogeneity' assumption
  \end{enumerate}
\end{enumerate}

SUTVA is violated when the control (unprotected) group is affected by
the treatment group. For biodiversity counts in MPA vs non-MPA
locations, marine life from the MPAs can travel to non-protected areas,
affecting the outcomes. Thus, violating the SUTVA assumption.
Additionally, MPAs can vary in their restrictiveness. Some MPAs are
no-take zones, while others have restrictions on seasons or gear,
violating SUTVA's `No hidden variation' parameter.

Exogeneity assumption refer to the outcome error being unbiased.
However, in a MPA treatment effect estimator, the error can be caused to
due to MPA placement (MPA location is chosen for a specific reason - not
randomly selected) or ecological covariates (depth and habitat).
Therefore, both SUTVA and exogeneity assumptions can be violated when
studying MPA effectiveness.

\begin{center}\rule{0.5\linewidth}{0.5pt}\end{center}

\section{EXTRA CREDIT}\label{extra-credit}

\begin{quote}
Use the recent lobster abundance data with observations collected up
until 2024 (\texttt{extracredit\_sblobstrs24.csv}) to run an analysis
evaluating the effect of MPA status on lobster counts using the same
focal variables.
\end{quote}

\begin{enumerate}
\def\labelenumi{\alph{enumi}.}
\tightlist
\item
  Create a new script for the analysis on the updated data
\item
  Run at least 3 regression models \& assess model diagnostics
\item
  Compare and contrast results with the analysis from the 2012-2018 data
  sample (\textasciitilde{} 2 paragraphs)
\end{enumerate}

\begin{center}\rule{0.5\linewidth}{0.5pt}\end{center}

\pandocbounded{\includegraphics[keepaspectratio]{figures/spiny1.png}}

\end{document}
